\chapter{The Agent Loop, State, and Memory}
\label{ch:agent-loop}

An agent loop is three steps repeated: think, act, observe. Everything difficult
about operating one comes from the fact that the observations accumulate. Each
tool result joins the context, stays there, is re-read on every subsequent turn,
and is paid for again each time.

This chapter is about the loop as a thing you run in production: how it
terminates, where state lives now that the protocol holds none of it, and how to
stop the context filling up with the remains of steps that finished long ago.

\section{The loop, instrumented}

Four bands per iteration, and if you cannot attribute a millisecond to one of
them you cannot optimise it.

\begin{py}
@dataclass
class Iteration:
    """One trip round the loop, with the milliseconds attributed."""
    index: int
    model_ms: float = 0.0
    transport_ms: float = 0.0
    consent_ms: float = 0.0
    tool_calls: int = 0
    parallel: bool = False
    input_tokens: int = 0
    output_tokens: int = 0
    cost_usd: float = 0.0
\end{py}

The loop itself:

\begin{py}
for step in range(self.max_steps):
    iteration = Iteration(index=step)

    # --- think
    turn = self.model.think("\n".join(context), catalogue)
    iteration.model_ms = turn.latency_ms
    iteration.input_tokens = turn.input_tokens
    iteration.output_tokens = turn.output_tokens
    iteration.cost_usd = turn.cost_usd

    if not turn.tool_calls:
        result.answer = turn.text
        result.stopped_because = "completed"
        break

    # --- act
    iteration.parallel = self.parallel_fanout and len(turn.tool_calls) > 1
    started = time.perf_counter()
    if iteration.parallel:
        results = self.host.call_tools_parallel(turn.tool_calls)
    else:
        results = [self.host.call_tool(c["name"], c.get("arguments"))
                   for c in turn.tool_calls]
    iteration.transport_ms = (time.perf_counter() - started) * 1000.0

    # --- observe
    for call, payload in zip(turn.tool_calls, results):
        context.append(f"[{call['name']}] {_summarise(payload)}")
\end{py}

Think, act, observe. The whole thing.

\section{Termination}

An agent loop that cannot stop is not an agent, it is an outage with a
personality. There are four ways out and a production loop needs all four.

\begin{table}[htbp]
\centering
\small
\begin{tabularx}{\textwidth}{@{}lXl@{}}
\toprule
\thead{Condition} & \thead{Means} & \thead{Meridian default} \\
\midrule
No tool calls    & The model is finished. The one you hope for & \\
Step budget      & Looping without progress & 8 \\
Token budget     & Context filling with tool output & task-dependent \\
Cost budget      & Somebody wired this to cron and went on holiday & task-dependent \\
\bottomrule
\end{tabularx}
\caption{Four exits. The last three are guardrails, and the third and fourth are
the ones people leave unset.}
\label{tab:termination}
\end{table}

\begin{py}
if self.token_budget is not None and result.total_tokens > self.token_budget:
    result.stopped_because = "token budget exhausted"
    result.answer = "Stopped: the context budget for this task ran out."
    break
if self.cost_budget_usd is not None and result.cost_usd > self.cost_budget_usd:
    result.stopped_because = "cost budget exhausted"
    result.answer = "Stopped: the dollar budget for this task ran out."
    break
\end{py}

Note that a budget stop still produces an \emph{answer}. The user asked a
question and deserves to be told why they are not getting one, which an
exception thrown three frames up does not accomplish.

\begin{py}
def test_step_budget_stops_a_runaway(self):
    looping = [[{"name": "fraud.screen_account",
                 "arguments": {"accountId": "ACC-1000"}}]] * 50
    model = StubModel(plan=looping, simulate_latency=False)
    result = AgentLoop(host, model, max_steps=4).run("loop forever")
    self.assertEqual(result.stopped_because, "step budget exhausted")
    self.assertEqual(len(result.iterations), 4)
\end{py}

\keyidea{Set all four. The model cannot be trusted to stop, and the server
cannot see enough to make it. Termination is the host's job, and an unset budget
is not a default, it is infinity.}

\subsection{Choosing the numbers}

Not arbitrary, and derivable from Chapter~\ref{ch:primer}.

\textbf{Step budget}: look at your traces, take the 95th percentile of
iterations for successful tasks, and add two. Meridian's tasks finish in three,
so eight is generous.

\textbf{Token budget}: your context window, minus a reserve for the final answer,
minus your catalogue cost times your step budget. If that arithmetic comes out
negative, your catalogue is too big and Chapter~\ref{ch:tools} is where you
should be.

\textbf{Cost budget}: what one task is worth. A \$0.012 task with a \$0.50 budget
has forty times the headroom it needs and still stops before anything alarming
happens.

\section{Where state lives now}

The protocol holds none. So state lives in exactly three places, and knowing
which is which prevents most of the confusion.

\begin{table}[htbp]
\centering
\small
\begin{tabularx}{\textwidth}{@{}lXX@{}}
\toprule
\thead{Kind} & \thead{Where} & \thead{Lifetime} \\
\midrule
Conversation & Host memory & The conversation \\
Cross-call server state & A handle passed as a tool argument & Server's retention policy \\
Durable domain state & A database, exposed as resources or tools & Forever \\
\bottomrule
\end{tabularx}
\caption{Three homes. Nothing lives in the connection, because the connection is
no longer a thing that means anything.}
\label{tab:state-homes}
\end{table}

The middle row is the one that changed. Before 2026, a server could keep a
per-session dictionary. Now it mints a handle and the model carries it.

The interesting consequence: \textbf{the model is now part of your state
machine}. It has to remember to pass \texttt{basket\_id} to the next call. Which
means the handle needs to be visible in a tool result the model actually reads,
and the tool description needs to say what to do with it.

\begin{py}
# The model is responsible for carrying `basket_id` forward; the server
# stores the cart contents under that key and looks them up on each call.
\end{py}

If that makes you slightly nervous, good. It should. It is why
Chapter~\ref{ch:building-servers} insists on a clear expiry error: when the model
loses the handle, and it will, the failure has to be self-explanatory enough for
it to recover.

\section{Context management}

Now the hard part.

Every tool result enters the context and stays there. Ten calls returning 2\,KB
each is 20\,KB of JSON, re-processed on every subsequent turn, most of it
irrelevant by the time you get to turn eight.

\bookfig[0.94]{xkcd-context-growth}{Six iterations of one task. The input
line is everything you are paying to re-read.}{context-growth}

\subsection{The cost argument is the weaker one}

The obvious reason to manage context is money, and this chapter's cost table
makes that case. There is a second reason, and in practice it decides more
outcomes.

A model reading 60,000 tokens is not the same model reading 6,000. Attention is
finite and spread across everything present, so each additional token of
irrelevant material competes with the relevant ones. Two effects follow, and
both are worth designing around.

\textbf{Position matters.} Material at the very start and very end of a long
context is used more reliably than material in the middle. A tool result buried
at turn three of a twelve-turn conversation is in the worst possible place: still
paid for on every turn, and least likely to be acted on. This is why appending
is not a neutral operation, and why the summarisation strategies below beat
simply having a bigger window.

\textbf{Contradictions accumulate.} At turn one the account had a score of 55. At
turn seven, after a re-score, it has 61. Both are in the context, both look
authoritative, and nothing marks the first as superseded. The model now has to
guess which one you meant, and a plausible wrong guess is worse than a missing
value, because a missing value produces a question and a wrong guess produces a
confident answer.

So the goal is not a smaller context. It is a context where everything present is
current and relevant, which is a different objective and sometimes points the
other way: re-injecting a result you dropped can be the right call, and spending
tokens on an explicit ``the score was re-run; the current figure is 61'' is
usually worth more than the tokens it costs.

\keyidea{Growing your context window raises the ceiling on what fits. It does
not raise the ceiling on what the model reliably uses. Those are different
limits, and the second one is the one your task success rate is bounded by.}

\subsection{Truncate, at minimum}

Meridian's approach is blunt and bounded, which beats unbounded:

\begin{py}
def _summarise(payload: dict, limit: int = 600) -> str:
    """Put the tool result into the context without putting *all* of it there.

    Structured content is preferred over the text mirror, because it is smaller
    and the model does not have to parse prose to find a number. Truncation is
    blunt, and Chapter 13 replaces it with something better, but blunt and
    bounded beats unbounded.
    """
    if payload.get("structuredContent") is not None:
        text = json.dumps(payload["structuredContent"], separators=(",", ":"))
    else:
        text = " ".join(block.get("text", "")
                        for block in payload.get("content", [])
                        if block.get("type") == "text")
    if payload.get("isError"):
        text = "ERROR: " + text
    return text if len(text) <= limit else text[:limit] + "...[truncated]"
\end{py}

Preferring \texttt{structuredContent} is worth more than it looks. It is smaller
than the prose mirror and the model does not have to parse English to find a
number.

\subsection{Better strategies, roughly in order of effort}

\textbf{Summarise old turns.} After N iterations, replace the oldest tool results
with a one-line summary. Costs a model call. Best when tasks are long and early
results stop mattering.

\textbf{Retrieval-augmented re-injection.} Keep full results outside the context,
in a store, and re-inject only what the current step needs. More machinery, and
the right answer when results are large and reuse is unpredictable.

\textbf{Sliding window.} Keep the last N results, drop the rest. Cheap. Wrong
whenever step nine needs step one, which is more often than you would like.

\textbf{Resource links instead of content.} Chapter~\ref{ch:resources} covered
this: return a link, let the model fetch only if it needs to. Zero tokens for
results nobody reads.

\begin{notebox}{Measure against task success, not against token count}
Every one of these makes token count go down. That is not the goal. The goal is
tasks that still succeed.

An aggressive sliding window can cut context by 60\,\% and quietly drop your
success rate by 15\,\%, and you will only find out from the eval suite in
Chapter~\ref{ch:measuring}. Never tune context strategy without one.
\end{notebox}

\section{Cache staleness in a long-running agent}

A subtle failure mode that only shows up in autonomous systems.

An agent runs for twenty minutes. It cached \texttt{tools/list} at minute zero
with a five-minute TTL. At minute six the cache is stale, and the agent is
mid-task.

The freshness check handles it, but only if you actually check on access rather
than on a timer. And there is a second-order problem: what if the catalogue
\emph{changed} in a way that invalidates the plan the model already made?

Three defences, in increasing order of paranoia:

\textbf{Subscribe.} A \texttt{listChanged} notification invalidates immediately,
so the window is milliseconds instead of minutes.

\textbf{Treat unexpected errors as invalidation signals.} A \texttt{-32602}
saying ``unknown tool'' on a tool you know exists means your catalogue is stale.
Refetch and retry once. The specification
explicitly permits refetching before TTL expiry when you have reason to believe
the data changed.

\textbf{Re-plan on catalogue change.} If the tools available to a task change
mid-task, the plan may no longer be valid. Expensive, and worth it for long
autonomous runs.

\section{Cost guardrails}

\begin{py}
@property
def cost_usd(self) -> float:
    return sum(i.cost_usd for i in self.iterations)
\end{py}

Per iteration, from tokens and price:

\begin{py}
def estimate_cost_usd(input_tokens: int, output_tokens: int) -> float:
    return (input_tokens * INPUT_USD_PER_MTOK / 1_000_000.0
            + output_tokens * OUTPUT_USD_PER_MTOK / 1_000_000.0)
\end{py}

\begin{measurebox}{Where Meridian's \$0.0125 goes}
\begin{tabular}{@{}lrrr@{}}
\toprule
\thead{Iteration} & \thead{Input} & \thead{Output} & \thead{Cost} \\
\midrule
1 & 1{,}252 & 12 & \$0.003936 \\
2 & 1{,}318 & 19 & \$0.004239 \\
3 & 1{,}400 & 8  & \$0.004320 \\
\midrule
\textbf{Total} & \textbf{3{,}970} & \textbf{39} & \textbf{\$0.012495} \\
\bottomrule
\end{tabular}

\vspace{6pt}
\footnotesize
Note the input column climbing while the output column does not. That is the
context filling with tool results, and it is why iteration three costs more than
iteration one despite doing less.
\end{measurebox}

That table is the whole economics of agent loops in six numbers. Input grows
monotonically because everything accumulates. Output stays tiny because the model
emits a tool call, not an essay. So \textbf{your cost is dominated by re-reading
what you already read}, and the lever is context management.

\subsection{Model routing}

Not every step needs your best model.

\begin{table}[htbp]
\centering
\small
\begin{tabularx}{\textwidth}{@{}lXl@{}}
\toprule
\thead{Step} & \thead{Needs} & \thead{Model} \\
\midrule
Planning the task           & Real reasoning & Large \\
Picking one of three tools  & Classification & Small \\
Formatting a result         & Transcription & Small \\
Deciding whether to stop    & Judgement & Large \\
Summarising for context     & Compression & Small \\
\bottomrule
\end{tabularx}
\caption{Route by what the step actually requires. Glue steps are the majority
of turns in a long task, and they are the cheapest to move.}
\label{tab:model-routing}
\end{table}

The measurement that matters is the quality/cost frontier, and the honest way to find it is to run your eval suite at each mix. It
is common to move the majority of turns to a small model and lose nothing
measurable, and it is also common to lose a lot, and which one you are in is not
predictable from first principles.

\section{Memory across sessions}

The protocol is stateless. Your application probably should not be.

Meridian's risk agent keeps 30 days of interaction history for an account, and
none of it lives in MCP. It lives in a database, exposed as a resource:

\begin{plain}
meridian://accounts/{account_id}/history
\end{plain}

Which is the whole trick. \textbf{Persistence is a domain concern, not a
protocol concern.} The protocol moved to statelessness precisely so it would stop
having an opinion about your storage.

Three layers, and it is worth being explicit about which is which:

\begin{description}[leftmargin=1.4em, style=nextline, itemsep=3pt]
\item[Conversation memory] This task. Lives in the host, dies with the
conversation.
\item[Session memory] This user, recently. Lives in the host's own store.
Survives a page refresh.
\item[Domain memory] Facts about the world. Lives in your database, reachable as
resources and tools. Survives everything, and is the only layer another system
can also read.
\end{description}

The mistake is putting domain memory in conversation memory, where it is
expensive, ephemeral, and invisible to every other consumer.

\section{Streaming, because perceived latency is the latency}

The Meridian baseline is 1,366\,ms, 96\,\% of it model inference. You cannot make
that faster. You can make it \emph{feel} faster, which is the actual user experience and
therefore the thing worth optimising.

\begin{itemize}[leftmargin=1.6em, itemsep=3pt]
\item \textbf{Stream the model's tokens} as they arrive. Time to first token
becomes the number the user feels, and it is a third of the total.
\item \textbf{Show tool calls as they happen.} ``Checking fraud signals...''
tells the user the system is working and, incidentally, gives them a chance to
notice it is doing something wrong.
\item \textbf{Forward progress notifications.} A server emitting
\texttt{notifications/progress} is offering you free UX. Take it.
\item \textbf{Stream partial results.} If iteration one produced a usable
answer, show it while iteration two refines.
\end{itemize}

\begin{py}
if total > 20 and i % 10 == 0:
    ctx.progress(i, total, f"scored {i} of {total}")
\end{py}

Note the guard again. Progress on a three-item loop is noise on the wire and
flicker in the UI.

\section{Putting it together}

\begin{py}
def test_loop_terminates_and_attributes_time(self):
    host = wire_host()
    model = StubModel(plan=[
        [{"name": "risk.assess_account_risk",
          "arguments": {"accountId": "ACC-1002"}}],
        [{"name": "fraud.screen_account",
          "arguments": {"accountId": "ACC-1002"}}],
        [],
    ], simulate_latency=False)
    result = AgentLoop(host, model).run("Assess ACC-1002 for credit and fraud.")

    self.assertEqual(result.stopped_because, "completed")
    self.assertEqual(len(result.iterations), 3)
    self.assertEqual(result.round_trips, 2)
    self.assertGreater(result.model_ms, 0)
    self.assertGreater(result.total_tokens, 0)
\end{py}

Three iterations for two tool calls. The third exists solely so the model can
read the second result and decide it is done, and that third turn costs a full
TTFT and 1,400 input tokens.

\keyidea{The last iteration of every loop is pure overhead: a model turn that
produces no tool call and exists only to say ``I am finished''. It is the price
of letting the model decide when to stop, and it is usually worth paying. But it
means a two-call task costs three turns, and that is the arithmetic to use when
estimating.}

\section{What to remember}

\begin{itemize}[leftmargin=1.6em, itemsep=3pt]
\item Think, act, observe. Attribute every millisecond to one of four bands.
\item Four termination conditions, all four set to a number. Budget stops still
produce an answer.
\item State lives in host memory, in server handles passed as arguments, or in
your database. Never in the connection.
\item The model is now part of your state machine, so handles must be visible and
expiry errors must be self-explanatory.
\item Context management is the dominant cost lever, because input grows every
turn while output does not.
\item Tune context strategy against task success, never against token count
alone.
\item Treat an unexpected ``unknown tool'' as a cache invalidation signal.
\item Route glue steps to a small model, and verify the frontier with evals
rather than intuition.
\item Persistence is a domain concern. Expose it as resources.
\item Stream everything. Perceived latency is the latency.
\end{itemize}

Next: what happens when the thing calling your tools is itself an agent.
